\documentclass{article}

\textwidth=6.5in
\textheight=8.5in
\voffset=-54pt
\oddsidemargin=0in
\evensidemargin=0in
\setlength{\parskip}{8pt}
\setlength{\parindent}{0pt}

\usepackage{workbook}

\usepackage[sols]{handout}

\begin{document}

\begin{center}
  \large\textsc{Problem Set 35 - Modeling with Differential Equations}
  
      \pspace{12pt}
  \begin{problemonly}\large Name: \rule{5cm}{0.15mm}\\ \end{problemonly}
\end{center}
\begin{grading}
  Total: 25 points 
  \begin{multicols}{3}
    \begin{enumerate}[nosep]
    \item 2 points
    \item 4 points
    \item 8 points (2, 2, 2, 2) 
    \item 5 points (1, 1, 1, 2)
    \item 4 points
    \item 2 points
    \end{enumerate}
  \end{multicols}
\end{grading}

\begin{enumerate}
\item

  For help, you may wish to read pp. 982-988 in Gottlieb.

  \begin{problem}
    A large garbage dump sits on the outskirts of a city. Garbage is being
    deposited at the dump at a rate of 40 tons per month. Scavengers and
    salvagers frequent the dump and haul off refuse from the site. The
    rate at which garbage is being hauled off is proportional to the
    tonnage at the site. Let $G(t)$ be the number of tons of garbage in
    the dump at time $t$, where $t$ is measured in months.

    Write a differential equation whose solution is $G(t)$. (By this we
    mean that the differential equation should involve $\frac{dG}{dt}$.)
    The basic framework is
    \[ \text{rate of change of } G = \text{rate of increase of garbage} -
    \text{rate of decrease of garbage}.\]   
  \end{problem}

  \pspace{\vfill}

  \begin{solution}
    The {\it rate of increase} of garbage is 40 tons per month.  The {\it
      rate of decrease} is proportional to the amount of refuse on the
    site, which is simply $G$.  Therefore, the equation
    \[ \text{rate of change of } G = \text{rate of increase of garbage} -
    \text{rate of decrease of garbage} \] can be rewritten as $\boxed{
      \frac{dG}{dt} = 40 - kG}$.
  \end{solution} 

\item 

  George pulls a pizza out of his oven and sets it on the counter to cool
  a little.  Let $P(t)$ be the pizza's temperature in degrees Fahrenheit $t$ minutes after it's
  pulled from the oven.
  \begin{enumerate}
  \item
    \begin{problem}
      Newton's Law of Cooling says that the rate of cooling (or heating)
      of an object is proportional to the temperature difference between
      the object and its surroundings.  If the pizza is sitting in a
      $70\,^\circ $F kitchen and is $190\,^\circ$F when it comes out of the oven,
      write a differential equation that models this situation.  Is there
      an initial condition?
    \end{problem}
    \pspace{\vfill}

    \begin{solution}
      Newton's Law of Cooling says that $\frac{dP}{dt}$ is proportional to
      $P - 70$, so $\frac{dP}{dt} = k (P - 70)$.  (The constant $k$ should
      be negative because $\frac{dP}{dt}$ is negative and $P - 70$ is
      positive for our pizza.)

      We also have an initial condition, $P(0) = 190$. 
    \end{solution}

  \item
    \begin{problem}
      When the pizza is $190\,^\circ$F, it's cooling at an instantaneous rate
      of $20\,^\circ$F per minute.  Incorporate this information into your
      differential equation.
    \end{problem}
    \pspace{\vfill}

    \begin{solution}
      We are told that, when $P = 190$, $\frac{dP}{dt} = -20$.  Plugging
      this information into our differential equation, $-20 = k (190 -
      70)$, so $k = - \frac{1}{6}$.  Therefore, our differential equation
      becomes $\boxed{\frac{dP}{dt} = - \frac{1}{6} (P - 70)}$. 
    \end{solution}

  \end{enumerate}

  \pspace{\newpage}

\item 

  Elmer takes out a \$100,000 loan for a house. He pays
  money back at a rate of \$9,000 per year. The bank charges him interest
  at a nominal annual interest rate of 6\%, compounded continuously.

  \begin{enumerate}
  \item
    \begin{problem}
      Make a continuous model of his economic situation by writing a
      differential equation whose solution is $B(t)$, the balance Elmer
      owes the bank $t$ years after taking out the loan.
    \end{problem}

    \pspace{\stretch{0.5}}

    \begin{solution} % originally by Henry Shin, modified by Janet
      The rate of change of the balance owed is $\frac{dB}{dt}$, so 
      \[ \frac{dB}{dt} = (\text{rate at which balance increases}) -
      (\text{rate at which balance decreases}).\] The rate at which the
      balance increases (due to interest) is $0.06B$.  The rate at which
      the balance decreases (because of Elmer's payments) is $9000$
      dollars per year.  So, $\boxed{\frac{dB}{dt} = 0.06B - 9000}$.

      The fact that Elmer takes out a loan of \$100,000 gives us an
      initial condition $\boxed{B(0) = 100000}$.
    \end{solution}

  \item
    \begin{problem}
      According to your model, which of the following gives the amount of
      money that Elmer owes $t$ years after taking out the loan?
    \end{problem}

    \begin{multicols}{2}
      \begin{enumerate}[label=\protect\maybecircle{\roman*.}]
      \plainitem $B(t) = (100,000 - 9,000 t) e^{0.06 t}$
      \circleitem $B(t) = 150,000 - 50,000 e^{0.06 t}$
      \plainitem $B(t) = 100,000e^{0.06t} - 9,000t$
      \plainitem $B(t) = 150,000 e^{0.06t} - 50,000$
      \end{enumerate}
    \end{multicols}

    \begin{solution}
      We simply need to determine which of the functions satisfies both
      the differential equation $\frac{dB}{dt} = 0.06 B - 9000$ and the
      initial condition $B(0) = 100,000$.  Let's look at each one:
      \begin{enumerate}[label=\roman*.]
      \item If $B(t) = (100,000 - 9,000 t) e^{0.06 t}$, then
        $\frac{dB}{dt} = -9000 e^{0.06 t} + (100,000 - 9,000 t) \cdot 0.06
        e^{0.06 t}$, while $0.06 B - 9000 = 0.06 (100,000 - 9,000 t)
        e^{0.06 t} - 9000$.  These are not equal.
      \item If $B(t) = 150,000 - 50,000 e^{0.06 t}$, then $\frac{dB}{dt} =
        -50,000 (0.06 e^{0.06 t})$ while $0.06 B - 9000 = 0.06 \left(
          150,000 - 50,000 e^{0.06 t} \right) - 9000 = 0.06 \left( -
          50,000 e^{0.06 t} \right)$.  These are equal, so we at least
        have a solution of the differential equation.  We should also see
        if this function satisfies the initial condition: $B(0) = 150,000
        - 50,000 = 100,000$, so it does.  Therefore, this function must be
        the amount of money that Elmer owes $t$ years after taking out the
        loan.  
      \end{enumerate}      
    \end{solution}

    \pspace{\stretch{2}}

  \item
    \begin{problem}
      According to this model, how long does it take Elmer to pay off the
      loan?  Please give both an exact answer and a decimal approximation.
    \end{problem}
    \pspace{\vfill}

    \begin{solution}
      We found that Elmer's balance at time $t$ is $B(t) = 150,000 -
      50,000 e^{0.06 t}$; we'd like to figure out when this is 0, so we
      solve
      \begin{align*}
        150,000 - 50,000 e^{0.06 t} &= 0 \\
        150,000 &= 50,000 e^{0.06 t} \\
        3 &= e^{0.06 t} \\
        \ln 3 &= 0.06 t \\
        \frac{\ln 3}{0.06} &= t 
      \end{align*}
      So, it takes Elmer $\boxed{\frac{\ln 3}{0.06} \approx 18.3102048}$
      years to pay off the loan. 
    \end{solution}

  \item
    \begin{problem}
      How much money does Elmer end up paying the bank in all?  (How does
      this compare to the amount he borrowed?)
    \end{problem}
    \pspace{\vfill}

    \begin{solution}
      Elmer pays the bank \$9,000 per year for $\frac{\ln 3}{0.06} \approx
      18.3102048$ years, so he pays a total of $(\$9,000 \text{ per year})
      \left( \frac{\ln 3}{0.06} \text{ years}\right) \boxed{\approx
        \$164,791.84}$.  (This is quite a bit more than the \$100,000 he
      borrowed!)
    \end{solution}

  \end{enumerate}

\pspace{\newpage}

\item \label{diffeq-model-rumor} You are studying how a rumor spreads
  through a town of $N$ people.  Let $R(t)$ be the number of people who
  have heard the rumor at time $t$.  
  \begin{enumerate}
  \item Write differential equations expressing the following three
    models.
    \begin{enumerate}
    \item
      \begin{problem}
        The rumor spreads at a rate proportional to the number of people
        who have heard it.
      \end{problem}
      \pspace{\vfill}

      \begin{solution}
        The spreading rate is $\frac{dR}{dt}$, and the number of people
        who have heard the rumor is $R$.  So, this model says that
        $\frac{dR}{dt}$ is proportional to $R$; that is,
        $\boxed{\frac{dR}{dt} = kR}$ for some positive constant $k$.  
      \end{solution}

    \item
      \begin{problem}
        The rumor spreads at a rate proportional to the number of people
        who have not heard it yet.
      \end{problem}
      \pspace{\vfill}

      \begin{solution}
        The spreading rate is $\frac{dR}{dt}$, and the number of people
        who have not yet heard the rumor is $N - R$ ($N$ is the total
        population, and $R$ of those people have already heard it).  So,
        this model says that $\frac{dR}{dt}$ is proportional to $N - R$;
        that is, $\boxed{\frac{dR}{dt} = k(N - R)}$ for some constant $k$.
      \end{solution}

    \item
      \begin{problem}
        The rumor spreads at a rate proportional to the product of the
        number of the people who have heard it and the number of people
        who have not yet heard it.
      \end{problem}
      \pspace{\vfill}

      \begin{solution}
        In this model, $\frac{dR}{dt}$ is proportional to the product $R
        (N - R)$.  We can write that as $\boxed{\frac{dR}{dt} = k R (N -
          R)}$ for some constant $k$.
      \end{solution}

    \end{enumerate}

  \item
    \begin{problem}
      Assuming the town is completely cut off from the outside world,
      which model is the most realistic?  (Hint: What should the spreading
      rate be if nobody has heard the rumor?  What if everybody in the
      town has heard the rumor?)
    \end{problem}
    \pspace{\stretch{1}}

    \begin{solution}
      If nobody has heard the rumor, it should not spread; nobody is able
      to spread it because nobody knows it!

      If everybody has heard the rumor, it should not spread either; there
      is nobody new to spread it to.

      So, when $R = 0$ or $R = N$, $\frac{dR}{dt}$ should be 0.  This is
      only true of the \fbox{third model}. 
    \end{solution}

  \end{enumerate}

  \probonly{\textit{Note:} Epidemiologists use similar models to model the
    spread of diseases.}

\pspace{\newpage}

\item \label{gottlieb-31.2.16-parts}

  Money in a bank account is earning interest at a nominal rate of 4\% per
  year compounded continuously. Withdrawals are made at a rate of \$8000
  per year.  

  \begin{enumerate}

  \item \label{gottlieb-31.2.16a}
    \begin{problem}
      Write a differential equation modeling the situation.
    \end{problem}
    \pspace{\stretch{0.5}}

    \begin{solution}
      The rate of change of the amount of money in the account is
      $\frac{dM}{dt}$, and we can break this down as
      \begin{align*}
        \frac{dM}{dt} = {} & (\text{rate at which money increases due to
          interest}) \\
        & - (\text{rate at which money decreases due to withdrawals})
      \end{align*}
      The rate at which money increases due to interest is $0.04 M$, while
      the rate of withdrawal is \$8000 per year.  So, $\boxed{
        \frac{dM}{dt} = 0.04M - 8000}$.
    \end{solution}

  \item \label{gottlieb-31.2.16d}

    \begin{problem}
      Show that $M(t) = Ce^{0.04t} - 8000t$ is \emph{not} a solution to
      the differential equation you got in \ref{gottlieb-31.2.16a}.
      (Here, $C$ is a constant.)
    \end{problem}
    \pspace{\vfill}

    \begin{solution}
      We just need to check that $\frac{dM}{dt} \neq 0.04 M - 8000$.  If
      $M(t) = Ce^{0.04t} - 8000t$, then $\frac{dM}{dt} = 0.04 C e^{0.04 t}
      - 8000$ while $0.04 M - 8000 = 0.04 (C e^{0.04 t} - 8000 t) - 8000$,
      so we see that $\frac{dM}{dt} \neq 0.04 M - 8000$.
    \end{solution}

  \end{enumerate}

\item  \begin{problem}
    \textbf{Reflect Back.} Sergei is looking back at
    \ref{gottlieb-31.2.16d}, and he says, ``I don't understand why $M(t) =
    Ce^{0.04t} - 8000t$ is not a solution.  If the account starts with
    $C$ dollars, then after $t$ years, interest will make the amount of
    money go up to $C e^{0.04 t}$, and it will have lost $8000 t$
    dollars due to withdrawals.  So, the total amount of money in the
    account should be $C e^{0.04 t} - 8000 t$ dollars.''

    What's wrong with Sergei's reasoning?
  \end{problem}
  \pspace{\vfill}

  \begin{solution}
    Sergei assumes incorrectly that the entire \$$C$ the account started
    with earns interest for $t$ years.  Some of that money gets withdrawn
    from the account and therefore stops earning interest. 
  \end{solution}

\end{enumerate}

\psreflection{}
\end{document}